\documentclass[10pt,a4paper]{amsart}
\usepackage{amsmath,amssymb,amsthm,graphicx}
% For commutative diagrams we use Xy-pic
\usepackage[all]{xy}

% We use 2cell for 2-commutative diagrams.
\xyoption{2cell}
% Theorem environments.
%
\theoremstyle{plain}
\newtheorem{theorem}[subsection]{Theorem}
\newtheorem{proposition}[subsection]{Proposition}
\newtheorem{lemma}[subsection]{Lemma}

\theoremstyle{definition}
\newtheorem{definition}[subsection]{Definition}
\newtheorem{example}[subsection]{Example}
\newtheorem{exercise}[subsection]{Exercise}
\newtheorem{situation}[subsection]{Situation}

\theoremstyle{remark}
\newtheorem{remark}[subsection]{Remark}
\newtheorem{remarks}[subsection]{Remarks}

\numberwithin{equation}{subsection}

% Macros
%
\def\lim{\mathop{\mathrm{lim}}\nolimits}
\def\colim{\mathop{\mathrm{colim}}\nolimits}
\def\Spec{\mathop{\mathrm{Spec}}}
\def\Hom{\mathop{\mathrm{Hom}}\nolimits}
\def\Ext{\mathop{\mathrm{Ext}}\nolimits}
\def\SheafHom{\mathop{\mathcal{H}\!\mathit{om}}\nolimits}
\def\SheafExt{\mathop{\mathcal{E}\!\mathit{xt}}\nolimits}
\def\Sch{\mathit{Sch}}
\def\Mor{\mathop{\mathrm{Mor}}\nolimits}
\def\Ob{\mathop{\mathrm{Ob}}\nolimits}
\def\Sh{\mathop{\mathit{Sh}}\nolimits}
\def\NL{\mathop{N\!L}\nolimits}
\def\CH{\mathop{\mathrm{CH}}\nolimits}
\def\proetale{{pro\text{-}\acute{e}tale}}
\def\etale{{\acute{e}tale}}
\def\QCoh{\mathit{QCoh}}
\def\Ker{\mathop{\mathrm{Ker}}}
\def\Im{\mathop{\mathrm{Im}}}
\def\Coker{\mathop{\mathrm{Coker}}}
\def\Coim{\mathop{\mathrm{Coim}}}

% Boxtimes
%
\DeclareMathSymbol{\boxtimes}{\mathbin}{AMSa}{"02}

%
% Macros for moduli stacks/spaces
%
\def\QCohstack{\mathcal{QC}\!\mathit{oh}}
\def\Cohstack{\mathcal{C}\!\mathit{oh}}
\def\Spacesstack{\mathcal{S}\!\mathit{paces}}
\def\Quotfunctor{\mathrm{Quot}}
\def\Hilbfunctor{\mathrm{Hilb}}
\def\Curvesstack{\mathcal{C}\!\mathit{urves}}
\def\Polarizedstack{\mathcal{P}\!\mathit{olarized}}
\def\Complexesstack{\mathcal{C}\!\mathit{omplexes}}
% \Pic is the operator that assigns to X its picard group, usage \Pic(X)
% \Picardstack_{X/B} denotes the Picard stack of X over B
% \Picardfunctor_{X/B} denotes the Picard functor of X over B
\def\Pic{\mathop{\mathrm{Pic}}\nolimits}
\def\Picardstack{\mathcal{P}\!\mathit{ic}}
\def\Picardfunctor{\mathrm{Pic}}
\def\Deformationcategory{\mathcal{D}\!\mathit{ef}}
\usepackage{cleveref}
\crefname{originaldefinition}{original definition}{original definitions}
\Crefname{originaldefinition}{Original definition}{Original definitions}
\makeatletter
\newlabel{algebra-definition-perfect}{{030Y}{source}{}{sourcepointer-1}{}}
\newlabel{algebra-definition-perfect@cref}{{[originaldefinition][0][]030Y}{[1][0][]source}{}{}{}}
\makeatother
\crefname{originaldefinition}{original definition}{original definitions}
\Crefname{originaldefinition}{Original definition}{Original definitions}
\makeatletter
\newlabel{algebra-definition-perfection}{{046X}{source}{}{sourcepointer-2}{}}
\newlabel{algebra-definition-perfection@cref}{{[originaldefinition][0][]046X}{[1][0][]source}{}{}{}}
\makeatother
\crefname{originallemma}{original lemma}{original lemmas}
\Crefname{originallemma}{Original lemma}{Original lemmas}
\makeatletter
\newlabel{algebra-lemma-colimit-syntomic}{{07BV}{source}{}{sourcepointer-3}{}}
\newlabel{algebra-lemma-colimit-syntomic@cref}{{[originallemma][0][]07BV}{[1][0][]source}{}{}{}}
\makeatother
\crefname{originallemma}{original lemma}{original lemmas}
\Crefname{originallemma}{Original lemma}{Original lemmas}
\makeatletter
\newlabel{algebra-lemma-p-ring-map}{{0BRA}{source}{}{sourcepointer-4}{}}
\newlabel{algebra-lemma-p-ring-map@cref}{{[originallemma][0][]0BRA}{[1][0][]source}{}{}{}}
\makeatother
\crefname{originalreference}{original reference}{original references}
\Crefname{originalreference}{Original reference}{Original references}
\makeatletter
\newlabel{lemma-smooth-base-change-general}{{0F09}{source}{}{sourcepointer-5}{}}
\newlabel{lemma-smooth-base-change-general@cref}{{[originalreference][0][]0F09}{[1][0][]source}{}{}{}}
\makeatother
\crefname{originalreference}{original reference}{original references}
\Crefname{originalreference}{Original reference}{Original references}
\makeatletter
\newlabel{proposition-topological-invariance}{{03SI}{source}{}{sourcepointer-6}{}}
\newlabel{proposition-topological-invariance@cref}{{[originalreference][0][]03SI}{[1][0][]source}{}{}{}}
\makeatother
\title{Base change of a pushforward along an extension of fields with torsion coefficients invertible in the base}
\author{The Stacks Project Authors}
\date{stacks-project:etale-cohomology.tex:local-snapshot-20261011; GNU FDL 1.2 or later}
\begin{document}
\maketitle
\noindent\textit{Editorial note.} All delivered mathematical text is the author's own, copied from the pinned Stacks Project chapter file etale-cohomology.tex; no statement or proof step was rewritten, reordered or removed.\par\medskip
\noindent\textit{Editorial note.} The counted proof is the author's complete argument: the proof is written out in full: the separable case through the colimit-of-smooth-algebras description, and the general case through the perfect closures with topological invariance.\par\medskip
\noindent\textit{Editorial note.} The label lemma-base-change-field-extension names the counted statement, and the proof environment immediately following it in the source is delivered verbatim as its proof.\par\medskip
\noindent\textit{Editorial note.} The counted statement is carried by the author's own lemma environment in the chapter source.\par\medskip
\noindent\textit{Editorial note.} The statement and proof are copied from the pinned chapter file; the six cited inputs are kept as external pointers with their canonical Stacks tags.\par\medskip
\noindent\textit{Allowed standard prerequisites (editorial).}\begin{itemize}
\item Separable extensions are filtered colimits of smooth algebras. Algebra, the lemma presenting a separable field extension as a filtered colimit of smooth algebras. Supplies the reduction to the general smooth base change lemma in the separable case.
\item Perfect closures and purely inseparable extensions. Algebra, the definition of the perfection of a field and the lemma that an extension of perfect fields yields universal homeomorphisms, together with Same chapter's topological invariance of etale cohomology. Supplies the reduction of the general case to the separable case by passing to perfect closures.
\item The general smooth base change statement. Same chapter, the base change lemma for a directed limit of smooth schemes, applied to the separable extension. Supplies the isomorphism in the separable case.
\end{itemize}\medskip
\ifdefined\hypertarget\hypertarget{sourcepointer-1}{}\fi
\noindent\textit{Source reference.} Original definition 030Y: sources/stacks-chow/tags.txt, source lines 3965--3965. Prerequisite pointer invoked at proof line(s) 18062; it names a standard fact of a sibling chapter that the counted argument uses, and every step of the counted proof built on it is retained verbatim in the delivered body.\par\medskip
\ifdefined\hypertarget\hypertarget{sourcepointer-2}{}\fi
\noindent\textit{Source reference.} Original definition 046X: sources/stacks-chow/tags.txt, source lines 5488--5488. Prerequisite pointer invoked at proof line(s) 18049; it names a standard fact of a sibling chapter that the counted argument uses, and every step of the counted proof built on it is retained verbatim in the delivered body.\par\medskip
\ifdefined\hypertarget\hypertarget{sourcepointer-3}{}\fi
\noindent\textit{Source reference.} Original lemma 07BV: sources/stacks-chow/tags.txt, source lines 9355--9355. Prerequisite pointer invoked at proof line(s) 18045; it names a standard fact of a sibling chapter that the counted argument uses, and every step of the counted proof built on it is retained verbatim in the delivered body.\par\medskip
\ifdefined\hypertarget\hypertarget{sourcepointer-4}{}\fi
\noindent\textit{Source reference.} Original lemma 0BRA: sources/stacks-chow/tags.txt, source lines 14763--14763. Prerequisite pointer invoked at proof line(s) 18053; it names a standard fact of a sibling chapter that the counted argument uses, and every step of the counted proof built on it is retained verbatim in the delivered body.\par\medskip
\ifdefined\hypertarget\hypertarget{sourcepointer-5}{}\fi
\noindent\textit{Source reference.} Original reference 0F09: sources/stacks-etale-cohomology/pinned-original-source.tex, source lines 17975--17975. Prerequisite pointer invoked at proof line(s) 18047; it names a standard fact of the same chapter that the counted argument uses, and every step of the counted proof built on it is retained verbatim in the delivered body.\par\medskip
\ifdefined\hypertarget\hypertarget{sourcepointer-6}{}\fi
\noindent\textit{Source reference.} Original reference 03SI: sources/stacks-etale-cohomology/pinned-original-source.tex, source lines 6548--6548. Prerequisite pointer invoked at proof line(s) 18059; it names a standard fact of the same chapter that the counted argument uses, and every step of the counted proof built on it is retained verbatim in the delivered body.\par\medskip
\begin{lemma}
\label{lemma-base-change-field-extension}
Let $L/K$ be an extension of fields. Let $g : T \to S$ be a quasi-compact
and quasi-separated morphism of schemes over $K$. Denote
$g_L : T_L \to S_L$ the base change of $g$ to $\Spec(L)$.
Let $E \in D^+(T_\etale)$ have cohomology sheaves whose stalks
are torsion of orders invertible in $K$. Let $E_L$ be the
pullback of $E$ to $(T_L)_\etale$. Then
$Rg_{L, *}E_L$ is the pullback of $Rg_*E$ to $S_L$.
\end{lemma}

\begin{proof}
If $L/K$ is separable, then $L$ is a filtered colimit of smooth
$K$-algebras, see
Algebra, Lemma \ref{algebra-lemma-colimit-syntomic}.
Thus the lemma in this case follows immediately from
Lemma \ref{lemma-smooth-base-change-general}.
In the general case, let $K'$ and $L'$ be the perfect closures
(Algebra, Definition \ref{algebra-definition-perfection})
of $K$ and $L$. Then $\Spec(K') \to \Spec(K)$ and
$\Spec(L') \to \Spec(L)$ are universal homeomorphisms
as $K'/K$ and $L'/L$ are purely inseparable
(see Algebra, Lemma \ref{algebra-lemma-p-ring-map}).
Thus we have $(T_{K'})_\etale = T_\etale$,
$(S_{K'})_\etale = S_\etale$,
$(T_{L'})_\etale = (T_L)\etale$, and
$(S_{L'})_\etale = (S_L)_\etale$ by
the topological invariance of \'etale cohomology, see
Proposition \ref{proposition-topological-invariance}.
This reduces the lemma to the case of the field
extension $L'/K'$ which is separable (by definition of
perfect fields, see Algebra, Definition \ref{algebra-definition-perfect}).
\end{proof}

\end{document}
